\section*{Problemas}

\subsection*{Enunciados de los problemas}

% Recordar
\begin{problema}
	\label{prob:39c2251}
	Una muestra aleatoria de $n = 36$ estudiantes obtuvo una calificación promedio muestral de $\bar{X} = 78$ puntos en un examen de matemáticas. Suponga que la desviación estándar poblacional es conocida y exactamente igual a $\sigma = 12$ puntos.
	\begin{enumerate}
		\item Construya un intervalo de confianza del 95\% para la media poblacional $\mu$.
		\item Construya un intervalo de confianza del 99\% para la media poblacional $\mu$.
		\item ¿Qué intervalo posee una amplitud mayor? Explique el motivo teórico por el cual la amplitud varía en función del nivel de confianza.
	\end{enumerate}
\end{problema}

% Comprender
\begin{problema}
	\label{prob:8fd99c8}
	Un fabricante de focos LED desea estimar la vida útil promedio de su nueva línea de productos. Toma una muestra aleatoria de $n = 16$ focos y registra un promedio de $\bar{X} = 1,200$ horas con una desviación estándar muestral $s = 80$ horas.
	\begin{enumerate}
		\item ¿Debe utilizar la distribución normal $Z$ o la distribución $t$ de Student para calcular el intervalo de confianza? Argumente su respuesta.
		\item Construya un intervalo de confianza del 95\% para la verdadera vida útil promedio.
		\item Interprete formalmente el intervalo resultante en el contexto del problema.
	\end{enumerate}
\end{problema}

% Aplicar
\begin{problema}
	\label{prob:6af5954}
	Una firma de telecomunicaciones requiere determinar el tiempo promedio real de las llamadas telefónicas de sus abonados corporativos. A partir de una auditoría sobre $n = 100$ llamadas, encuentra una media muestral $\bar{X} = 4.5$ minutos con $s = 2.0$ minutos.
	\begin{enumerate}
		\item Construya el intervalo de confianza del 95\% para la duración media $\mu$.
		\item Determine la magnitud del margen de error de estimación ($E$) actual.
		\item Si la gerencia técnica requiere reducir el margen de error máximo tolerado a $E = 0.2$ minutos al mismo nivel del 95\% de confianza, deduzca y calcule el tamaño de muestra mínimo $n$ necesario.
	\end{enumerate}
\end{problema}

% Analizar
\begin{problema}
	\label{prob:0257f5f}
	\textbf{Deducción analítica de la cota conservadora del tamaño muestral para una proporción Bernoulli:} al diseñar una encuesta para estimar la verdadera proporción poblacional $p \in (0,1)$ con un margen de error máximo admitido $E$ y una confianza $(1-\alpha)$, la expresión algebraica que define el semiancho del intervalo asintótico normal es:
	\begin{align}
		E = Z_{\alpha/2} \sqrt{\frac{p(1-p)}{n}}.
	\end{align}
	\begin{enumerate}
		\item Demuestre rigurosamente, mediante cálculo diferencial, que la varianza muestral teórica de Bernoulli $g(p) = p(1-p)$ alcanza un máximo global absoluto en el dominio $p \in [0,1]$ exactamente en $p = 0.5$.
		\item Demuestre que si no se dispone de información a priori sobre $p$, la fórmula que garantiza que el margen de error nunca excederá el límite $E$ es $n = \left\lceil \dfrac{Z_{\alpha/2}^2}{4 E^2} \right\rceil$.
		\item Una consultora de ciencia de datos requiere estimar la intención de voto hacia una iniciativa pública de salud con un margen de error máximo del $3\%$ ($E = 0.03$) y un nivel de confianza del 95\%. Calcule el número mínimo de encuestas válidas que deberá garantizar.
	\end{enumerate}
\end{problema}

% Evaluar
\begin{problema}
	\label{prob:a3ef454}
	\textbf{Fundamentación epistemológica del intervalo de confianza:} considere la construcción formal del intervalo de confianza empírico $(1-\alpha)$ bajo normalidad, $[L(X), U(X)] = \left[ \bar{X} - Z_{\alpha/2}\frac{\sigma}{\sqrt{n}}, \bar{X} + Z_{\alpha/2}\frac{\sigma}{\sqrt{n}} \right]$. Suponga que, tras tomar la muestra real de 36 estudiantes del Problema \ref{prob:39c2251}, se publica el intervalo numérico $[74.08, 81.92]$.
	\begin{enumerate}
		\item Evalúe por qué, desde la axiomática estricta de la estadística frecuentista, es teóricamente erróneo enunciar que ``la probabilidad de que la verdadera calificación media $\mu$ esté dentro de $[74.08, 81.92]$ es del $95\%$''.
		\item Defina rigurosamente qué significa una ``confianza del 95\%'' desde la perspectiva del muestreo repetido infinito, y contrástelo con el concepto epistemológico del \textbf{intervalo de credibilidad Bayesiano}, donde sí es válido asignar una probabilidad directa al parámetro dados los datos.
	\end{enumerate}
\end{problema}

% Crear
\begin{problema}
	\label{prob:18d6b96}
	Diseñe un problema original que ilustre la dualidad entre intervalos de confianza y pruebas de hipótesis (distinto de calificaciones, focos o llamadas ya vistos): un laboratorio farmacéutico afirma que sus tabletas contienen exactamente $500$ mg de principio activo. Un control de calidad externo toma una muestra de $n=20$ tabletas, con media muestral $\bar{X}=496$ mg y desviación estándar muestral $s=9$ mg. Construya el intervalo de confianza del 95\% para el contenido medio real, determine si el valor afirmado ($500$ mg) está contenido en el intervalo, y explique la conclusión en términos de una prueba de hipótesis bilateral implícita al mismo nivel de significancia.
\end{problema}

\subsection*{Soluciones detalladas}

% Recordar
\begin{solproblema}[prob:39c2251]
	\begin{enumerate}
		\item Con $Z_{0.025}=1.96$ y $\text{SE}=12/\sqrt{36}=2$:
		\begin{align}
			IC_{95\%} = 78 \pm 1.96(2) = 78 \pm 3.92 = [74.08,\;81.92].
		\end{align}
		\item Con $Z_{0.005}=2.576$:
		\begin{align}
			IC_{99\%} = 78 \pm 2.576(2) = 78 \pm 5.152 = [72.85,\;83.15].
		\end{align}
		\item Las amplitudes son $A_{95\%}=7.84$ y $A_{99\%}=10.30$ puntos. El intervalo del 99\% es más amplio: para incrementar la confianza de no dejar fuera al verdadero $\mu$ manteniendo constante la información muestral ($n$), es necesario expandir el margen de captura, abarcando una fracción mayor del área bajo las colas de la normal.
	\end{enumerate}
\end{solproblema}

% Comprender
\begin{solproblema}[prob:8fd99c8]
	\begin{enumerate}
		\item Debe usarse la \textbf{distribución $t$ de Student}: el tamaño de muestra es pequeño ($n=16<30$, lo que impide invocar el TLC) y la desviación estándar poblacional $\sigma$ es desconocida (se sustituye por $s=80$, introduciendo un error de estimación adicional que exige las colas más pesadas de $t$).
		\item Con $\nu=15$ y $t_{0.025,15}=2.131$:
		\begin{align}
			IC_{95\%} = 1{,}200 \pm 2.131\left(\frac{80}{\sqrt{16}}\right) = 1{,}200 \pm 2.131(20) = 1{,}200 \pm 42.62 = [1{,}157.38,\;1{,}242.62] \text{ horas.}
		\end{align}
		\item Si repitiéramos el proceso de muestreo (extraer 16 focos y calcular el intervalo) indefinidamente, en el 95\% de esas muestras la vida útil media verdadera $\mu$ de toda la producción quedaría capturada entre 1,157.38 y 1,242.62 horas.
	\end{enumerate}
\end{solproblema}

% Aplicar
\begin{solproblema}[prob:6af5954]
	\begin{enumerate}
		\item Con $n=100\ge30$ (TLC aplicable) y $Z_{0.025}=1.96$:
		\begin{align}
			IC_{95\%} = 4.5 \pm 1.96\left(\frac{2.0}{\sqrt{100}}\right) = 4.5 \pm 1.96(0.2) = 4.5 \pm 0.392 = [4.108,\;4.892] \text{ min.}
		\end{align}
		\item El margen de error actual es $E=0.392$ minutos.
		\item Partiendo de $E=Z_{\alpha/2}\dfrac{s}{\sqrt{n}}$, despejamos $n=\left(\dfrac{Z_{\alpha/2}\cdot s}{E}\right)^2$. Sustituyendo:
		\begin{align}
			n = \left(\frac{1.96\times2.0}{0.2}\right)^2 = (19.6)^2 = 384.16 \implies n_{\text{óptimo}} = \lceil384.16\rceil = 385 \text{ llamadas.}
		\end{align}
	\end{enumerate}
\end{solproblema}

% Analizar
\begin{solproblema}[prob:0257f5f]
	\begin{enumerate}
		\item Sea $g(p)=p(1-p)=p-p^2$. Derivando: $g'(p)=1-2p=0 \implies p^*=0.5$. Como $g''(p)=-2<0$, es un máximo global, con valor $g(0.5)=0.25$.
		\item Partiendo de $Z_{\alpha/2}\sqrt{p(1-p)/n}\le E$, elevamos al cuadrado: $n\ge Z_{\alpha/2}^2\,p(1-p)/E^2$. Como $p$ es desconocida, sustituimos su cota superior $p(1-p)\le0.25$ (Parte 1) para que la desigualdad valga para cualquier $p$:
		\begin{align}
			n \ge \frac{Z_{\alpha/2}^2(0.25)}{E^2} = \frac{Z_{\alpha/2}^2}{4E^2} \implies n = \left\lceil\frac{Z_{\alpha/2}^2}{4E^2}\right\rceil.
		\end{align}
		\item Con $Z_{0.025}=1.96$ y $E=0.03$:
		\begin{align}
			n = \frac{(1.96)^2}{4(0.03)^2} = \frac{3.8416}{0.0036} \approx 1{,}067.11 \implies n_{\text{óptimo}} = 1{,}068 \text{ encuestas.}
		\end{align}
	\end{enumerate}
\end{solproblema}

% Evaluar
\begin{solproblema}[prob:a3ef454]
	\begin{enumerate}
		\item En el paradigma frecuentista, el parámetro poblacional $\mu$ es una \textbf{constante fija de la naturaleza}, no una variable aleatoria. Una vez extraída la muestra, $\bar{X}$ colapsa al valor observado $\bar{x}=78$, y el intervalo $[74.08,81.92]$ pasa a ser un conjunto numérico completamente determinista. El evento $\mu\in[74.08,81.92]$ ya no tiene ninguna aleatoriedad asociada: o el verdadero $\mu$ está en ese rango (probabilidad $1$) o no lo está (probabilidad $0$). Asignarle el número $0.95$ confunde la probabilidad del \emph{procedimiento} de muestreo con la probabilidad de un parámetro fijo.
		\item \textbf{Interpretación frecuentista correcta:} la confianza del 95\% es una propiedad del \emph{proceso generador del intervalo}, no del intervalo numérico particular. Si se repitiera el muestreo aleatorio un número infinito de veces, generando miles de intervalos $[L(X_k),U(X_k)]$, exactamente el 95\% de esos intervalos aleatorios envolverían al verdadero $\mu$. \textbf{Contraste bayesiano:} en la inferencia Bayesiana, el parámetro $\theta$ sí se trata como una variable aleatoria con distribución a priori $\pi(\theta)$; al condicionar en los datos observados, el Teorema de Bayes produce una distribución a posteriori $\pi(\theta\mid X=x)$, a partir de la cual el intervalo de credibilidad $[A,B]$ satisface $P(A\le\theta\le B\mid X=x)=0.95$ de forma directa y legítima —justo la afirmación que el enfoque frecuentista no puede hacer sobre un intervalo ya calculado.
	\end{enumerate}
\end{solproblema}

% Crear
\begin{solproblema}[prob:18d6b96]
	Con $n=20$ ($\nu=19$), $t_{0.025,19}=2.093$, y $\text{SE}=s/\sqrt{n}=9/\sqrt{20}\approx2.0125$:
	\begin{align}
		IC_{95\%} = 496 \pm 2.093(2.0125) \approx 496 \pm 4.212 = [491.79,\;500.21] \text{ mg.}
	\end{align}
	El valor afirmado ($500$ mg) \textbf{sí está contenido} en el intervalo, ya que $491.79 < 500 < 500.21$. Por la dualidad entre estimación por intervalos y pruebas de hipótesis, si se realizara una prueba $H_0:\mu=500$ vs. $H_a:\mu\neq500$ al nivel $\alpha=0.05$, \textbf{no existiría evidencia suficiente para rechazar $H_0$}: aunque la media muestral ($496$ mg) está por debajo de la especificación, la discrepancia observada es admisible como fluctuación aleatoria del muestreo, y el laboratorio de control de calidad no puede refutar concluyentemente la afirmación del fabricante.
\end{solproblema}
